% New English translation from checked cote 25, site 14 (PDF page 15).
% AI-translated by GPT-6 Astra (OpenAI, ChatGPT Pro mode).
\scanpage{25/14}
\mspage{\hand{11}}

this, let Z be the set of points at which X is not smooth over
k; its inverse image \var{\inserted[typed]{$Z_{\hand{\eta}}$ in $Y_{\hand{\eta}}$}}{$Z_{\hand{\eta}}$ in $Y_{\hand{\eta}}$}{insertion}{ega25_25_014-v0001} is therefore (by 2.12.) the set of
points of $Y_{\hand{\eta}}$ at which $Y_{\hand{\eta}}$ is not smooth over k($\hand{\eta}$). Now the codimension
of Z in X at x equals that of $Z_{\hand{\eta}}$ in $Y_{\hand{\eta}}$ at y by
flatness (ref par.6). Thus one is \hand{$>$}k iff the other is,
which completes the proof.
\uline{Corollary} 2.15. Let Z be the set of points of X at which
X is not smooth over k (resp. is not geom $R_k$ over k, resp. is
\var{\margin{\unsure{//}}}{\placedmargin{\unsure{//}}}{marginal note placed}{ega25_25_014-v0002} not separable over \unsure{R|k}, resp. is not \var{\inserted[typed]{geom}}{geom}{insertion}{ega25_25_014-v0003} normal over k). For $Y_{\hand{\eta}}$ to be
smooth (resp. geom $R_k$, resp separable, resp geom normal) over k($\hand{\eta}$ ), it
is n and s that Z be finite.

\var{\margin{$|$}}{\placedmargin{$|$}}{marginal note placed}{ega25_25_014-v0004} From the standpoint of exposition, 2.14. and 2.15. contain 2.1\unsure{2|3}. and
2.13. (which we could therefore dispense with stating separately); on the
other hand, the corollary is more important in practice than the theorem,
which could be given as a proposition or preliminary lemma, so as to
glorify the corollary instead.

We can vary \var{\xout{\unsure{xxxxbyxi}}}{}{typewriter strike-out}{ega25_25_014-v0005} the case of property (iii):
\uline{Corollary} 2.16. (Still assuming f an immersion \var{\inserted[typed]{and F moreover coherent}}{and F moreover coherent}{insertion}{ega25_25_014-v0006}), under the
\var{\margin{\unsure{//}}}{\placedmargin{\unsure{//}}}{marginal note placed}{ega25_25_014-v0007} conditions of 2.7., for $Z_i$ to be nonembedded, it is n and s that $Z_i$ be so.
When this is the case, the radicial mu\hand{l}tiplicity of F at $Z_i$ \var{\xout{(ref)}}{}{typewriter strike-out}{ega25_25_014-v0008}\var{\deleted{rel}}{}{deletion}{ega25_25_014-v0009} \var{\inserted{rel.}}{rel.}{insertion}{ega25_25_014-v0010}
to k (ref) equals that of $G_{\hand{\eta}}$ at $Z_i$ rel to k($\hand{\eta}$).

The first assertion is contained in the last assertion of
2.7. For the second, since $Y_{\hand{\eta}} \hand{\longrightarrow} X$ is flat and hence the functor F \hand{$\rightsquigarrow$} $G_{\hand{\eta}}$
exact, we are reduced, by restriction to a neighborhood of the gen pt of $Z_i$
and dévissage, to the case F = $\underline{O}_{Z_i}$ \var{\inserted[typed]{and we may assume $Z_i = X$.}}{and we may assume $Z_i = X$.}{insertion}{ega25_25_014-v0011} . Moreover, we can begin by
reducing to the case k alg closed, \var{\inserted[typed]{then X is separable over k.}}{then X is separable over k.}{insertion}{ega25_25_014-v0012} Then \var{\xout{\unsure{$\underline{O}$xxxxx$\underline{O}_{Z_i}$xxxxx}}}{}{typewriter strike-out}{ega25_25_014-v0013} the assertion is
contained in 2.15. (iii). We conclude as usual:
\tnote{25-014-1}{Nonembedded renders non immergé, in the terminology of components used here. Both occurrences of Z with subscript i in the criterion in 2.16 are retained. The source has a stray b in n'estb and splits d'hab itude; these spellings remain in the French. The source's abbreviations geom, rel, gen pt and alg closed stand for geometrically, relative, generic point and algebraically closed.}

\var{\margin{14}}{\placedmargin{14}}{marginal note placed}{ega25_25_014-v0014}
